\section{Frequency estimates and both wave variance directions}
\label{S:wave-variance}

We now apply the covariance and temporal estimates of
Section~\ref{S:masked-covariance} to the full wave phase. We retain
Hypothesis~\ref{H:homog}, write $a = \alpha_0$, and choose and fix
\begin{equation}\label{E:variance-gamma-range}
  0 < \gamma < \min\left(\frac12,\frac{2-\alpha}{2}\right).
\end{equation}
The interval is nonempty because $\alpha+a < 2$. We use
$\langle\xi\rangle = \sqrt{1+|\xi|^2}$ and set
\begin{equation}\label{E:variance-scale-profile}
  B_\alpha(h) =
  \begin{cases}
    1, & \alpha < 1,\\
    1+\log(1/h), & \alpha = 1,\\
    h^{1-\alpha}, & \alpha > 1,
  \end{cases}
  \qquad 0 < h \le 1.
\end{equation}

\subsection{A scalar row estimate for a measure}

Let $\nu$ be a positive Radon measure satisfying
\begin{equation}\label{E:frequency-translated-growth}
  \nu(B(y,r)) \le C r^\alpha,
  \qquad y \in \R^d,\quad r > 0,
\end{equation}
where $C \ge 0$. If the estimate is initially known for open balls, apply
it at radius $r+\epsilon$ and let $\epsilon \downarrow 0$ to obtain the
closed-ball estimate. Neither evenness nor a density is assumed here.
Using~\eqref{E:masked-frequency-pair}, define
\begin{align}
  \Phi_h(\xi,\eta)
  &:= \frac{\mathcal Q_h(|\xi|,|\eta|)}
  {\langle\xi\rangle\langle\eta\rangle},
  \label{E:frequency-row-kernel}\\
  J_h(\xi)
  &:= \langle\xi\rangle^{2\gamma}
  \int_{\R^d}\Phi_h(\xi,\xi-\lambda)\nu(\ud\lambda).
  \label{E:frequency-row-integral}
\end{align}
These are initially nonnegative extended integrals. In particular, the
notation denotes a pushforward of $\nu$, not integration against an
assumed spectral density.

\begin{proposition}[Uniform scalar row]\label{P:frequency-row}
  There is a finite $K = K(\alpha,a,\gamma)$, independent of
  $d,\nu,C,h$ and $\xi$, such that
  \begin{equation}\label{E:frequency-row-bound}
    \sup_{\xi\in\R^d}J_h(\xi) \le K C B_\alpha(h),
    \qquad 0 < h \le 1.
  \end{equation}
\end{proposition}

\begin{proof}
  If $C = 0$, a countable covering by finite balls gives $\nu = 0$, and
  the result is immediate. Suppose $C > 0$. For each $\xi$, let $\nu_\xi$
  be the pushforward under $\lambda \mapsto \eta = \xi-\lambda$.
  Then $\nu_\xi(B(0,R)) \le CR^\alpha$. For $q \ge 1$, put
  \[
    A_\xi(q) = \int_{|\eta|\le5q}\langle\eta\rangle^{-1}
    \nu_\xi(\ud\eta).
  \]
  Choose $N = \lceil\log_2(5q)\rceil$, so $5q \le 2^N < 10q$.
  Splitting into the unit ball and the annuli
  $2^j < |\eta| \le 2^{j+1}$ gives
  \[
    A_\xi(q) \le C\left(1+2^\alpha
    \sum_{j=0}^{N-1}2^{j(\alpha-1)}\right).
  \]
  The geometric sum is bounded when $\alpha < 1$, is $N$ when
  $\alpha = 1$, and is at most
  $2^{N(\alpha-1)}/(2^{\alpha-1}-1)$ when $\alpha > 1$. Hence,
  \begin{equation}\label{E:frequency-inverse-bessel}
    A_\xi(q) \le K_\alpha C
    \begin{cases}
      1, & \alpha < 1,\\
      1+\log q, & \alpha = 1,\\
      q^{\alpha-1}, & \alpha > 1,
    \end{cases}
  \end{equation}
  with a constant depending only on $\alpha$.

  Set $q = \max(1,|\xi|)$ and $\epsilon = 1-2\gamma > 0$. On
  $|\lambda| \le 4q$ we have $|\eta| \le 5q$. Both bracket powers in
  $\mathcal Q_h$ are at most one, since $a-1 < 0$. The contribution of
  this region to $J_h(\xi)$ is at most
  \[
    2\langle\xi\rangle^{-\epsilon}A_\xi(q)
    \le 2q^{-\epsilon}A_\xi(q).
  \]
  If $\alpha < 1$, this is uniformly bounded by~\eqref{E:frequency-inverse-bessel}.
  If $\alpha = 1$, the bound follows from
  \[
    e^{-\epsilon u}(1+u) \le 1+\frac1{e\epsilon},
    \qquad u = \log q \ge 0,
  \]
  since $u e^{-\epsilon u}$ has maximum $1/(e\epsilon)$. If
  $\alpha > 1$, the remaining power is $q^{\alpha+2\gamma-2}$, whose
  exponent is negative by~\eqref{E:variance-gamma-range}. Thus, the near
  contribution is at most a constant times $C$ in all three cases.

  On the far region write $r = |\lambda| > 4q$. Then
  \[
    |\eta| \ge \frac{3r}4,\qquad
    \big||\eta|-|\xi|\big| \ge \frac r2,\qquad
    |\eta|+|\xi| \ge \frac r2.
  \]
  Each oscillatory bracket is at least $\langle hr\rangle/2$, whereas
  $\langle\eta\rangle^{-1} \le 4/(3r)$ and
  $\langle\xi\rangle^{-\epsilon} \le 1$. Consequently, the far
  contribution is at most
  \[
    \frac43\,2^{2-a}\int_{|\lambda|>4q}
    |\lambda|^{-1}\langle h|\lambda|\rangle^{a-1}\nu(\ud\lambda).
  \]
  Enlarge the domain to $|\lambda| > 1$ and split it into dyadic annuli.
  Applying~\eqref{E:frequency-translated-growth} to their masses bounds
  this expression by a constant depending only on $\alpha,a$ times
  $C S(h)$, where
  \[
    S(h) = \sum_{j=0}^\infty
    2^{j(\alpha-1)}\langle h2^j\rangle^{a-1}.
  \]
  Choose $L = \lfloor\log_2(1/h)\rfloor \ge 0$. For $j \le L$ the
  bracket power is at most one. For $j \ge L+1$ it is at most
  $(h2^j)^{a-1}$. Since $\alpha+a-2 < 0$, the tail is
  \[
    \frac{h^{a-1}2^{(L+1)(\alpha+a-2)}}{1-2^{\alpha+a-2}}
    \le \frac{h^{1-\alpha}}{1-2^{\alpha+a-2}}.
  \]
  The inequality uses $2^{L+1} > 1/h$ and the negative exponent.
  The finite sum is bounded by $1/(1-2^{\alpha-1})$ for $\alpha < 1$,
  by $1+\log_2(1/h)$ for $\alpha = 1$, and by
  \[
    \frac{2^{\alpha-1}h^{1-\alpha}}{2^{\alpha-1}-1}
  \]
  for $\alpha > 1$. Since $h^{1-\alpha} \le 1$ in the first case,
  we obtain $S(h) \le K_{\alpha,a}B_\alpha(h)$. Combining the near and
  far bounds, and using $B_\alpha(h) \ge 1$, proves the assertion and
  finiteness of the original integral. All constants arose from numerical
  shell estimates and are independent of $d$.
  This proves Proposition~\ref{P:frequency-row}.
\end{proof}

\begin{corollary}[A time exponent above one half]\label{C:frequency-time}
  For every
  \begin{equation}\label{E:variance-beta-range}
    \frac12 < \beta < \min\left(1-\frac a2,
    \frac{3-a-\alpha}{2}\right),
  \end{equation}
  there is a finite $K_\beta = K_\beta(\alpha,a,\gamma,\beta)$ such that
  \begin{equation}\label{E:frequency-time-bound}
    h^{2-a}\sup_\xi J_h(\xi) \le K_\beta C h^{2\beta},
    \qquad 0 < h \le 1.
  \end{equation}
\end{corollary}

\begin{proof}
  Both upper endpoints in~\eqref{E:variance-beta-range} exceed one half,
  by $a < 1$ and $\alpha+a < 2$. If $\alpha < 1$, use
  $2-a > 2\beta$ in Proposition~\ref{P:frequency-row}. If $\alpha > 1$,
  use $3-a-\alpha > 2\beta$. For $\alpha = 1$, set
  $\delta = 2-a-2\beta > 0$ and $u = \log(1/h)$. Then
  \[
    h^\delta(1+\log(1/h))
    = e^{-\delta u}(1+u) \le 1+\frac1{e\delta}.
  \]
  This absorbs the critical logarithm. No endpoint value of $\beta$ is used.
  This proves Corollary~\ref{C:frequency-time}.
\end{proof}

\subsection{Actual forms on the full phase}

Choose and fix $T > 0$. Set $\Lambda = (1-\Delta)^{1/2}$,
$s = 1/2$ and $E_r = H^r(\R^d)\times H^{r-1}(\R^d)$. On the
complexification of $E = E_s$, use inner products linear in the first slot.
The map
\[
  Jz(\xi) = \left(\langle\xi\rangle^{1/2}\widehat z_1(\xi),
  \langle\xi\rangle^{-1/2}\widehat z_2(\xi)\right)
\]
is unitary onto $L^2(\R^d;\mathbb C^2)$. The full massless wave group has
Fourier matrix
\begin{equation}\label{E:variance-free-matrix}
  S(t,\xi) =
  \begin{pmatrix}
    \cos(t|\xi|) & \sin(t|\xi|)/|\xi|\\
    -|\xi|\sin(t|\xi|) & \cos(t|\xi|)
  \end{pmatrix},
\end{equation}
where the sine quotient is $t$ at zero. Let $B(b)(u,v) = (0,bu)$.
Choose and fix a real even nonnegative $\zeta \in C_c^\infty(\R^d)$
equal to one on $B_1$ and zero off $B_2$, and put $\zeta_R(x) = \zeta(x/R)$
for $R > 0$. For $I = [t_0,t_0+h] \subset [0,T]$, $0 < h \le 1$, and
Schwartz phases $z,w$, put
\begin{align}
  v_z(t) &= [S(t)z]_1,
  \label{E:variance-displacement}\\
  G_w(t) &= -\frac{\sin(t|D|)}{|D|}\Lambda w_1
  +\cos(t|D|)\Lambda^{-1}w_2,
  \label{E:variance-dual-wave}\\
  F_{z,w}(t,x) &= 1_I(t)v_z(t,x)\overline{G_w(t,x)},
  \qquad \ell_R(I;z,w) = W_{\mathbb C}(\zeta_RF_{z,w}),
  \label{E:variance-actual-form}
\end{align}
where $W_{\mathbb C}(F) = W(\operatorname{Re}F)+iW(\operatorname{Im}F)$.
There is no observation cutoff in this definition.

We first assume $\Gamma \ne 0$. If $\Gamma = 0$, all actual forms and
the column maps below are zero, so their bounds hold directly.
For the nonzero case, let $\nu_R$ be~\eqref{E:masked-auxiliary-measure} and set
\begin{equation}\label{E:variance-actual-growth-constant}
  L_\zeta = \|\widehat\zeta\|_1,\qquad
  C_\nu = c_d^2L_\zeta^2C_\Gamma.
\end{equation}
Proposition~\ref{P:masked-auxiliary-norm} gives a positive even tempered
Radon measure with $\nu_R(B(y,r)) \le C_\nu r^\alpha$, uniformly in $R$.

\begin{proposition}[Actual form and phase kernel]\label{P:actual-wave-form}
  The tests $\zeta_RF_{z,w}$ in~\eqref{E:variance-actual-form} belong to the
  original noise completion. The form $\ell_R$ is linear in $z$ and
  conjugate linear in $w$, and uses the same $W$ as the model.
  For a smooth deterministic coefficient $b$ with
  common compact spatial support and bounded spatial derivatives on $[0,T]$,
  \begin{equation}\label{E:variance-physical-adjoint}
    \langle S(-t)B(b(t))S(t)z,w\rangle_E
    = \int b(t,x)v_z(t,x)\overline{G_w(t,x)}\ud x.
  \end{equation}
  For a real compact smooth $\chi$, multiplication $D_\chi$ on both phase
  entries has adjoint
  \begin{equation}\label{E:variance-cutoff-adjoint}
    D_\chi^{*,E}w
    = \left(\Lambda^{-1}\chi\Lambda w_1,
    \Lambda\chi\Lambda^{-1}w_2\right).
  \end{equation}
  With $U = Jz$, $V = Jw$, $r = |\xi|$, $\rho = |\eta|$ and the
  continuous functions $c_r,b_r$ of Lemma~\ref{L:masked-continuous-wave}, set
  \begin{equation}\label{E:variance-normalized-kernel}
    \mathcal K(t;\xi,\eta)
    = \frac1{\sqrt{\langle\xi\rangle\langle\eta\rangle}}
    \begin{pmatrix}
      -b_rc_\rho & -b_rb_\rho\\
      c_rc_\rho & c_rb_\rho
    \end{pmatrix}(t).
  \end{equation}
  Then
  \begin{equation}\label{E:variance-normalized-covariance}
    \E|\ell_R(I;z,w)|^2
    \le \int\left\|\int\sum_{i,j=1}^2
    1_I\mathcal K_{ij}(\cdot;\eta+\lambda,\eta)
    U_j(\eta)\overline{V_i(\eta+\lambda)}\ud\eta
    \right\|_{H_a^{\mathbb C}}^2\nu_R(\ud\lambda).
  \end{equation}
\end{proposition}

\begin{proof}
  The wave multipliers are smooth at zero through the continuous sine
  quotient, with polynomially bounded frequency derivatives on finite time
  intervals. Bessel multipliers preserve Schwartz space. Thus, $v_z,G_w$
  and their product are continuous in all Schwartz seminorms on $[0,T]$.
  For tests supported in a fixed compact $K$, the absolute spatial pairing
  is bounded by
  \[
    |K|\Gamma(K-K)\|F(t)\|_\infty\|H(u)\|_\infty.
  \]
  Indeed, only $y \in K-K$ contributes to
  $\int\Gamma(\ud y)\int F(t,x)\overline{H(u,x-y)}\ud x$.
  Combine this with the temporal Schur bound and approximate the interval
  indicator and the continuous time dependence by smooth interior temporal
  tests, keeping the spatial support of $\zeta_R$. This admits the masked
  product in the original completion in every allowed dimension.

  We also identify its auxiliary and spectral limits. Approximate a
  Schwartz-valued function on $I$ by time partitions with frozen spatial
  values, and then replace the finitely many indicators by compact smooth
  temporal functions. A diagonal sequence $F_n$ of finite temporal tensors
  converges in $L^2$ time in every Schwartz seminorm. Fourier integration
  by parts and the spatial weighted supremum bounds give, for sufficiently
  large fixed $N$,
  \[
    \|\widehat{F_n-F}(\xi)\|_{L^2(0,T)}
    \le \epsilon_n\langle\xi\rangle^{-N},\qquad \epsilon_n \to 0.
  \]
  Choose $N$ so that the square weight is integrable against both tempered
  measures $\mu$ and $\nu_R$. Lemma~\ref{L:masked-time-completion} gives
  convergence in their spectral $H_a$ norms. Multiplication by the fixed
  $\zeta_R$ preserves these $L^2$ Schwartz-seminorm limits, so the masked
  differences converge in the $\mu$ spectral norm too. The preceding
  compact-support estimate also gives convergence in the actual raw
  completion. Hence, the auxiliary extension and original-noise limit
  agree on these products. Constants for this admission argument may
  depend on $R$; they are not the uniform variance constants below.

  For real tensor tests $f,g$, the Hermitian symmetry of their spatial
  transforms gives
  \[
    A(-\xi) = \overline{A(\xi)},\qquad
    A(\xi) = \langle\widehat f(\xi),\widehat g(\xi)\rangle_{H_a^{\mathbb C}}.
  \]
  Integration against either even $\mu$ or even $\nu_R$ cancels the
  imaginary mixed term. Expanding the squared norm of
  $\widehat f+i\widehat g$, then using the preceding approximation, extends
  Proposition~\ref{P:masked-noise-completion} to the present complex tests:
  \begin{equation}\label{E:variance-complex-noise}
    \E|W_{\mathbb C}(\zeta_RF)|^2
    = (2\pi)^d\int\|\widehat{\zeta_RF}(\xi)\|_{H_a^{\mathbb C}}^2\mu(\ud\xi)
    \le (2\pi)^d\int\|\widehat F(\lambda)\|_{H_a^{\mathbb C}}^2\nu_R(\ud\lambda).
  \end{equation}
  This is a bound for the actual masked noise, not a declaration that its
  law is stationary. Complex linearity of $W_{\mathbb C}$ and linearity of
  $v_z,G_w$ prove the stated sesquilinearity.

  Smooth multiplication is bounded on each Sobolev space: Fourier
  convolution, the inequality
  $\langle\xi\rangle^r\langle\eta\rangle^{-r}
  \le 2^{|r|/2}\langle\xi-\eta\rangle^{|r|}$ and weighted Schwartz decay
  reduce this to the $L^1*L^2$ bound. The free matrix is strongly continuous
  and bounded on $E$ on compact time intervals, directly by its weighted
  Fourier entries. Thus, the ordinary first-response integral exists on
  states. Since
  $S(-t)(0,f)=(-\sin(t|D|)f/|D|,\cos(t|D|)f)$, moving the real
  self-adjoint multipliers and the two Sobolev weights onto $w$ proves
  \eqref{E:variance-physical-adjoint}. The same weighted duality gives
  \eqref{E:variance-cutoff-adjoint}. Inserting this adjoint in $G_w$ is
  the correctly observed form; $D_\chi$ need not be self-adjoint on $E$.
  This deterministic identity does not identify a positive-width source
  with the raw noise or assert a source limit.

  The definitions give exactly
  \begin{align*}
    \widehat v_z(t,\eta)
    &= \langle\eta\rangle^{-1/2}
    [c_\rho(t)U_1(\eta)+b_\rho(t)U_2(\eta)],\\
    \widehat G_w(t,\xi)
    &= \langle\xi\rangle^{-1/2}
    [-b_r(t)V_1(\xi)+c_r(t)V_2(\xi)].
  \end{align*}
  Their product gives~\eqref{E:variance-normalized-kernel}. At zero
  frequencies one uses the continuous sine quotient, not division by zero;
  at $\xi=\eta=0$ the matrix is
  $\left(\begin{smallmatrix}-t&-t^2\\1&t\end{smallmatrix}\right)$.
  Fourier multiplication with the conjugate yields
  \[
    \widehat F_{z,w}(t,-\lambda)
    = c_d1_I(t)\int\sum_{i,j=1}^2\mathcal K_{ij}(t;\eta+\lambda,\eta)
    U_j(\eta)\overline{V_i(\eta+\lambda)}\ud\eta.
  \]
  The spatial integral is absolutely convergent uniformly on the finite
  time interval by Schwartz decay. The Schur bound also makes it a
  Bochner $H_a^{\mathbb C}$ integral. Reflect $\lambda$ using evenness of
  $\nu_R$ in~\eqref{E:variance-complex-noise}. Since $(2\pi)^dc_d^2=1$,
  this proves~\eqref{E:variance-normalized-covariance} with no remaining
  Fourier factor.
  This proves Proposition~\ref{P:actual-wave-form}.
\end{proof}

\subsection{Bessel's inequality and the two variances}

\begin{lemma}[Hilbert-valued finite Bessel inequality]\label{L:phase-bessel}
  Let $H$ be a separable complex Hilbert space, let $\{u_k\}$ be a finite
  orthonormal family in $L^2(\R^d;\mathbb C^2)$, and let
  $g=(g_1,g_2) \in L^2(\R^d;H^2)$. Then
  \begin{equation}\label{E:phase-bessel}
    \sum_k\left\|\int\sum_{j=1}^2u_{k,j}(\eta)g_j(\eta)\ud\eta\right\|_H^2
    \le \int\sum_{j=1}^2\|g_j(\eta)\|_H^2\ud\eta.
  \end{equation}
  The same bound holds with $\overline{u_{k,j}}$ in place of $u_{k,j}$.
\end{lemma}

\begin{proof}
  Expand $g$ in a countable orthonormal basis of $H$. On each scalar
  coefficient apply scalar Bessel to the orthonormal family
  $\{\overline{u_k}\}$, or to $\{u_k\}$ for the conjugated version.
  Scalar finite Bessel follows by subtracting an orthogonal finite
  projection and expanding its nonnegative squared norm. Sum over the
  $H$ basis and use Parseval and nonnegative Tonelli. The Bochner integrals
  exist by $L^2$ Cauchy--Schwarz. A finite-dimensional $H$ uses its finite
  basis, and the zero space is immediate. No independence is involved.
  This proves Lemma~\ref{L:phase-bessel}.
\end{proof}

\begin{proposition}[Both actual variance directions]\label{P:wave-two-variances}
  There is a finite $C_T$, depending on the fixed model, mask,
  $a,\alpha,\gamma,T$, but not on $R,I$ or the orthonormal family, such
  that for every finite complex $E$-orthonormal Schwartz family $\{e_k\}$
  and every Schwartz phase $z$,
  \begin{align}
    \sum_k\E|\ell_R(I;e_k,z)|^2
    &\le C_T h^{2-a}B_\alpha(h)\|z\|_{E_{s-\gamma}}^2,
    \label{E:wave-input-variance}\\
    \sum_k\E|\ell_R(I;z,e_k)|^2
    &\le C_T h^{2-a}B_\alpha(h)\|z\|_{E_{s-\gamma}}^2.
    \label{E:wave-output-variance}
  \end{align}
  For every $\beta$ in~\eqref{E:variance-beta-range}, either right side
  may be replaced by $C_{T,\beta}h^{2\beta}\|z\|_{E_{s-\gamma}}^2$.
\end{proposition}

\begin{proof}
  Each entry in~\eqref{E:variance-normalized-kernel} is a product of two
  continuous wave factors times
  $(\langle\xi\rangle\langle\eta\rangle)^{-1/2}$, with a possible minus
  sign. Lemma~\ref{L:masked-continuous-wave} therefore gives
  \begin{equation}\label{E:variance-kernel-time-bound}
    \|1_I\mathcal K_{ij}(\cdot;\xi,\eta)\|_{H_a^{\mathbb C}}^2
    \le A_T h^{2-a}\Phi_h(\xi,\eta),
    \qquad A_T = C_1(a)(1+T)^4.
  \end{equation}
  The positive even Radon measure $\nu_R$ satisfies the actual growth bound
  \eqref{E:variance-actual-growth-constant}. Thus, all hypotheses of
  Proposition~\ref{P:frequency-row} hold, and
  \begin{equation}\label{E:variance-actual-row-bound}
    \int\Phi_h(\xi,\xi-\lambda)\nu_R(\ud\lambda)
    \le K C_\nu B_\alpha(h)\langle\xi\rangle^{-2\gamma}.
  \end{equation}
  No nonnegative physical inverse covariance of $\nu_R$ is needed.

  The kernels used in Lemma~\ref{L:phase-bessel} belong to its stated
  $L^2$ space for each fixed $\lambda$. Indeed,
  $|c_r(t)| \le 1$ and $|b_r(t)| \le \sqrt{2}\max(1,T)$ on $[0,T]$,
  by the high- and low-frequency representations. The kernels are
  continuous into $L^2(0,T)$ in their frequency variables, hence into
  $H_a^{\mathbb C}$ by Schur. Separability gives strong measurability.
  Multiplication by shifted Schwartz components of $U$ or $V$ then gives
  the required $L^2$ integrability. Later $\lambda$ integrals are of
  nonnegative functions; their finiteness follows from the estimates below.

  First fix $V = Jz$ and let $U_k = Je_k$ be the input family. In
  \eqref{E:variance-normalized-covariance} set
  \[
    g_j^\lambda(\eta) = \sum_{i=1}^2
    1_I\mathcal K_{ij}(\cdot;\eta+\lambda,\eta)
    \overline{V_i(\eta+\lambda)}.
  \]
  Lemma~\ref{L:phase-bessel} and Cauchy--Schwarz over the two output
  components give
  \begin{align*}
    \sum_j\|g_j^\lambda(\eta)\|_{H_a^{\mathbb C}}^2
    &\le 2\sum_{i,j}|V_i(\eta+\lambda)|^2
    \|1_I\mathcal K_{ij}(\cdot;\eta+\lambda,\eta)\|_{H_a^{\mathbb C}}^2\\
    &\le 4A_T h^{2-a}\Phi_h(\eta+\lambda,\eta)
    \sum_i|V_i(\eta+\lambda)|^2.
  \end{align*}
  Change variables $\xi = \eta+\lambda$ in the Lebesgue integral for
  each $\lambda$, apply Tonelli and~\eqref{E:variance-actual-row-bound},
  and use $\|z\|_{E_{s-\gamma}} = \|\langle\xi\rangle^{-\gamma}Jz\|_2$.
  This proves~\eqref{E:wave-input-variance} with $C_T = 4A_TKC_\nu$.

  Next fix $U = Jz$ and let $V_k = Je_k$ be the output family. After
  $\eta = \xi-\lambda$ in~\eqref{E:variance-normalized-covariance}, put
  \[
    f_i^\lambda(\xi) = \sum_{j=1}^2
    1_I\mathcal K_{ij}(\cdot;\xi,\xi-\lambda)U_j(\xi-\lambda).
  \]
  Apply the conjugated version of Lemma~\ref{L:phase-bessel}.
  Cauchy--Schwarz over the two input components bounds its right side by
  \[
    4A_T h^{2-a}\int\Phi_h(\xi,\xi-\lambda)
    \sum_j|U_j(\xi-\lambda)|^2\ud\xi.
  \]
  After the change of variables and Tonelli, the row at fixed $\eta$ is
  $\int\Phi_h(\eta+\lambda,\eta)\nu_R(\ud\lambda)$. Reflect $\lambda$
  using the actual evenness of $\nu_R$ and then use symmetry of $\Phi_h$.
  This row becomes
  $\int\Phi_h(\eta,\eta-\lambda)\nu_R(\ud\lambda)$, which is bounded by
  \eqref{E:variance-actual-row-bound} with $\xi = \eta$. This proves
  \eqref{E:wave-output-variance} with the same $C_T$.

  Use Corollary~\ref{C:frequency-time} for $h^{2-a}$ times the row integral
  at the two applications above to obtain the asserted $h^{2\beta}$ bounds,
  including strict absorption of the critical logarithm. The
  constants are uniform in $R,I$ and family size. The $R$-dependent
  admission estimate was not used as this uniform variance bound.
  This proves Proposition~\ref{P:wave-two-variances}.
\end{proof}

\subsection{Same-noise mean-square columns}

Write $E_r^{\mathrm{real}}$ for the real phase space in dimension $d$.

\begin{proposition}[Column and adjoint-column maps]\label{P:wave-square-columns}
  On the same probability space there are deterministic bounded real
  linear maps
  \[
    m_{R,I},n_{R,I}:E_{s-\gamma}^{\mathrm{real}}\longrightarrow
    L^2(\Omega;E_s^{\mathrm{real}}),
  \]
  with squared bounds $C_T h^{2-a}B_\alpha(h)$, such that, for every
  fixed real Schwartz pair $z,w$,
  \begin{equation}\label{E:variance-square-columns}
    \langle m_{R,I}z,w\rangle_E
    = \ell_R(I;z,w) = \langle z,n_{R,I}w\rangle_E
    \quad\text{in }L^2(\Omega).
  \end{equation}
  Fixed-vector values and finite families of values are centered Gaussian.
\end{proposition}

\begin{proof}
  Choose a countable real Schwartz orthonormal basis $\{e_k\}$ of $E_s$.
  Such a basis follows by Gram--Schmidt from a countable dense Schwartz
  family, discarding zero residuals. Schwartz density follows by smooth
  Fourier cutoff approximation in the weighted $L^2$ norms:
  approximate by compact smooth Fourier coordinates, replace $g(\xi)$
  by $(g(\xi)+\overline{g(-\xi)})/2$, and apply the inverse Bessel/Fourier
  maps. The resulting phases are real Schwartz, and Gram--Schmidt
  preserves that class. The same
  density holds in $E_{s-\gamma}$.

  For fixed real Schwartz $z$, define
  \[
    m_N(z) = \sum_{k\le N}\ell_R(I;z,e_k)e_k,
    \qquad n_N(z) = \sum_{k\le N}\ell_R(I;e_k,z)e_k.
  \]
  Both are finite Gaussian sums on the original $W$. By
  Proposition~\ref{P:wave-two-variances}, their coefficient expected squares
  have finite sum. Orthogonality gives the squared tail norms as the
  corresponding sums of expected squares. Thus, both sequences converge
  in $L^2(\Omega;E_s)$ to $m(z),n(z)$ with the asserted bounds, without
  independence of the coefficients. Finite-sum measurability and separable
  $E_s$ give strong measurability of the limits. The bounded real linear
  maps extend uniquely from Schwartz phases to $E_{s-\gamma}$.

  If $w_N$ is the $E_s$ projection of a Schwartz $w$ onto the first $N$
  basis vectors, then
  $\langle m_N(z),w\rangle_E = \ell_R(I;z,w_N)$. The singleton form of
  \eqref{E:wave-output-variance}, after normalizing $w-w_N$ in $E_s$,
  gives
  \[
    \E|\ell_R(I;z,w-w_N)|^2
    \le C_T h^{2-a}B_\alpha(h)\|z\|_{E_{s-\gamma}}^2\|w-w_N\|_{E_s}^2.
  \]
  The zero residual is immediate. The right side tends to zero, giving
  the first identity in~\eqref{E:variance-square-columns}. For the second,
  project $z$ instead and use the singleton input bound
  \eqref{E:wave-input-variance}. A fixed countable Schwartz orthonormal
  basis and Parseval show basis independence of the constructed limits.

  Continuous linear functionals of the limits, and finite linear
  combinations of such functionals, are $L^2$ limits of centered Gaussian
  variables with convergent variances. Their characteristic functions
  therefore converge to those of centered Gaussians. The same reasoning
  after the bounded input extension proves joint Gaussianity for arbitrary
  fixed inputs and finite families of the two maps.
  This proves Proposition~\ref{P:wave-square-columns}.
\end{proof}

These are maps into $L^2(\Omega;E_s)$ for fixed intervals. No common
uncountable almost-sure version, random bounded operator, Hilbert--Schmidt
realization, time path, source convergence or new solution-moment bound
is inferred by this section alone. Both variance directions are retained
in Section~\ref{S:actual-operator-driver}.
